\section{MPPI Controller Tuning}
\label{appendix:mppi_tuning}

\noindent\textit{This appendix covers MPPI parameter tuning, implementation features,
and compute performance. The world model used by MPPI is described in
Appendix~\ref{appendix:world_model}. Per-trial MPPI results are in
Appendix~\ref{appendix:per_trial}.}

\subsection{Algorithm Overview}

Model Predictive Path Integral (MPPI) is a sampling-based MPC
algorithm~\cite{hewing2020learning,williams2017mppi}. At each control step it
samples $N$ random action sequences around the previous solution, rolls each
one out through the learned world model on CPU (at each of the $H$ sequential
LSTM steps the $N$ candidate rollouts are advanced together as one batch,
i.e.\ parallel across samples but sequential along the horizon), and scores
the resulting
trajectories by cost. It then forms an
importance-weighted mean of the sampled sequences, applies the first action of
that averaged plan, and shifts the trajectory left to warm-start the next
step.

\subsection{Cost Function}

\begin{equation}\label{eq:mppi_cost}
J^{(i)} = \sum_{t=1}^{H} \left[ \mathbf{x}_t^T Q \mathbf{x}_t + R u_t^2 + c_{\text{wall}}(x_t) \right] + \mathbf{x}_H^T P \mathbf{x}_H
\end{equation}

where:
\begin{itemize}
  \item $Q = \text{diag}([0, 0.1, 30, 0.1])$: heavy ball-angle penalty
  \item $R = 0.001$: very low control effort penalty
  \item $c_{\text{wall}}(x_t) = 100 \cdot \max(0, |x^{\text{cart}}_t| - x_{\text{limit}})$: wall penalty on the cart position $x^{\text{cart}}_t$ (first coordinate of $\mathbf{x}_t$)
  \item $P = \text{diag}([1, 2, 100, 5])$: terminal cost (LQR approximation), disabled in the deployed configuration
\end{itemize}

\subsection{Importance Weighting}

\begin{equation}\label{eq:mppi_weight}
w^{(i)} = \frac{\exp(-(J^{(i)} - J_{\min}) / \lambda)}{\sum_j \exp(-(J^{(j)} - J_{\min}) / \lambda)}
\end{equation}

\begin{equation}\label{eq:mppi_weighted_mean}
\bar{U}^* = \sum_{i=1}^{N} w^{(i)} U^{(i)}
\end{equation}

The temperature parameter $\lambda$ controls how greedily the weights concentrate
on low-cost trajectories: small $\lambda$ makes the distribution sharply
peaked on the best trajectory; large $\lambda$ averages more uniformly.
We use $\lambda=0.4$ to balance greedy exploitation with smooth averaging.

\subsection{Action Sampling with Temporal Smoothing}

\begin{equation}\label{eq:mppi_smoothing}
u_t^{(i)} = \beta \cdot (\bar{u}_t + \epsilon_t^{(i)}) + (1-\beta) \cdot u_{t-1}^{(i)}, \quad \beta = 0.9
\end{equation}
with the base case $u_0^{(i)} = \bar{u}_0 + \epsilon_0^{(i)}$ (the first step is
unsmoothed) and the recursion above for $t \ge 1$.

Temporal smoothing ($\beta=0.9$) correlates noise across consecutive
timesteps, producing smoother candidate trajectories that are more
physically realizable on a velocity-controlled actuator and reduce
high-frequency chatter in the selected action sequence.

\subsection{Final Tuned Parameters}

Table~\ref{tab:mppi_params} lists the parameters used for every reported MPPI
run; the coordinate sweep that selected them is described next.

\begin{table}[!ht]
\centering
\caption{MPPI tuned parameters.}
\label{tab:mppi_params}
\begin{tabular}{lll}
\toprule
\textbf{Parameter} & \textbf{Value} & \textbf{Description} \\
\midrule
$Q$ & $[0, 0.1, 30, 0.1]$ & State cost weights \\
$R$ & 0.001 & Control effort penalty \\
$\lambda$ & 0.4 & Temperature (lower = more greedy) \\
Horizon ($H$) & 30 & Planning steps (1.5\,s lookahead) \\
Samples ($N$) & 800 & Parallel rollouts \\
$T_s$ & 0.05\,s & Control timestep (20\,Hz) \\
\bottomrule
\end{tabular}
\end{table}

\subsection{Tuning Method}

We tune MPPI by a coordinate sweep in simulation, varying each of the ball-angle weight $Q[2]$,
the control penalty $R$, the temperature $\lambda$, and the horizon $H$ in turn about a
fixed center ($Q[2]{=}30, H{=}30, R{=}0.001, \lambda{=}0.4$) at $N=400$ rollouts, scoring each configuration by the number of initial
conditions settled and the mean settling time across 7 difficulty-stratified
ICs (30\,s per trial). We run the sweep on the released MPPI world model
(Appendix~\ref{appendix:wm_provenance}), trained on the dataset
\texttt{arcball\_cont\_1M\_calib196\_160.h5}; because it is conducted in
simulation and is robust to minor model updates, the resulting parameter
values transfer to the final hardware checkpoint.

\subsection{Tuning Sweep Results}

The sweep varied one factor at a time and evaluated $Q[2] \in \{15,20,25,30,40,50\}$,
$H \in \{20,25,30,35\}$, $R \in \{0.001,0.005,0.01\}$,
$\lambda \in \{0.3,0.4,0.5\}$ at $N=400$ rollouts
with 7 difficulty-stratified ICs (30\,s per trial) in simulation.
An earlier parameterization ($Q[2]=80$, $R=0.01$, $H=15$, $N=800$) is
included as a baseline reference point.

\begin{table}[!t]
\centering
\caption{MPPI tuning sweep: selected configurations.
  All-settled (7/7) runs ranked by mean settling time;
  partial-success runs shown below.}
\begin{tabular}{lcccccc}
\toprule
\textbf{Config} & $Q[2]$ & $R$ & $\lambda$ & $H$ & \textbf{Settled} & \textbf{Mean\,(s)} \\
\midrule
\multicolumn{7}{l}{\textit{All 7 ICs settled}} \\
early config & 80 & 0.01 & 0.4 & 15 & 7/7 & 4.89 \\
q30\_h25 & 30 & 0.001 & 0.4 & 25 & 7/7 & 6.74 \\
\textbf{q30\_h30} & \textbf{30} & \textbf{0.001} & \textbf{0.4} & \textbf{30} & \textbf{7/7} & \textbf{7.16} \\
q30\_h35 & 30 & 0.001 & 0.4 & 35 & 7/7 & 7.37 \\
q30\_h20 & 30 & 0.001 & 0.4 & 20 & 7/7 & 8.64 \\
q15\_h30 & 15 & 0.001 & 0.4 & 30 & 7/7 & 9.61 \\
q50\_h30 & 50 & 0.001 & 0.4 & 30 & 7/7 & 9.65 \\
\midrule
\multicolumn{7}{l}{\textit{Partial success}} \\
q20\_h30 & 20 & 0.001 & 0.4 & 30 & 6/7 & 5.74 \\
q30\_lam05 & 30 & 0.001 & 0.5 & 30 & 6/7 & 7.46 \\
q30\_r005 & 30 & 0.005 & 0.4 & 30 & 5/7 & 6.95 \\
q30\_r01 & 30 & 0.01 & 0.4 & 30 & 5/7 & 8.43 \\
\bottomrule
\end{tabular}
\end{table}

The hardware evaluation uses $Q[2]=30$, $R=0.001$, $\lambda=0.4$, $H=30$,
$N=800$ (bolded row above, with the sample count raised to 800 for
hardware; even at this higher count the mean compute stays within the
50\,ms deadline, see Table~\ref{tab:mppi_runtime}). This configuration
achieves 7/7 in simulation with $H=30$ providing sufficient lookahead;
increasing $H$ beyond 30 does not improve success rate and raises
compute cost. Several sweep configurations reach 7/7 with lower mean settling than the
deployed one (including the earlier $Q[2]=80$, $R=0.01$, $H=15$ config and
$Q[2]=25$, $H=30$), so the deployed settings are not the settling-optimal
choice by the tabulated score. We nonetheless deploy $Q[2]=30$, $R=0.001$,
$H=30$ because its higher $Q[2]$ and lower $R$ give smoother, less reactive
planning; since the sweep is run in simulation and the settling differences
among 7/7 configurations are small, this is a qualitative engineering choice
rather than a settling-optimal or measured-hardware selection.

\subsection{Advanced Implementation Features}

The deployed MPPI controller implements several features beyond the basic
sampling loop (Table~\ref{tab:mppi_features}); the table lists the full set
available in the code, and features marked $^\dagger$ are disabled for the
final evaluation. In the released configuration used for the 50-trial
evaluation, context warm-up and warm-start are active, while model error
correction and the terminal cost are turned off and adaptive noise has no
effect because the deployed sampler draws uniform noise (the $\sigma$
schedule applies only to Gaussian sampling). The cost function above gives
the full form; its terminal term $\mathbf{x}_H^T P \mathbf{x}_H$ is therefore
inactive in the deployed run.

\begin{table}[!ht]
\centering
\caption{MPPI implementation features. Features marked with $^\dagger$ exist
  in code but are disabled for the final evaluation (best performance
  achieved without them).}
\label{tab:mppi_features}
\footnotesize
\setlength{\tabcolsep}{3pt}
\begin{tabular}{lp{0.66\columnwidth}}
\toprule
\textbf{Feature} & \textbf{Description} \\
\midrule
Context warm-up & Last $K=10$ observations fed through LSTM before planning \\
Model error correction$^\dagger$ & EMA of prediction error applied as bias \\
Terminal cost$^\dagger$ & LQR cost-to-go extends effective planning horizon \\
Adaptive noise$^\dagger$ & $\sigma$ decays (0.95$\times$) when good; increases when struggling (inactive under uniform sampling) \\
Warm-start & Previous solution shifted left; last action repeated \\
Wall override$^\dagger$ & Hard override when cart near boundary (disabled) \\
Cart clamping$^\dagger$ & Physical wall simulated during rollouts (disabled) \\
\bottomrule
\end{tabular}
\end{table}

\subsection{Compute Performance}

MPPI is the only controller that approaches the real-time control
deadline (Table~\ref{tab:mppi_runtime}); the classical and on-policy
controllers run in under 1\,ms, and the remaining controllers
(IQL, NMPC, PPO~(WM), and the off-policy RL policies) sit in the
2--8\,ms range, still far below the 50\,ms deadline
(Figure~\ref{fig:compute_time}).
Although the world model was trained on a GPU, MPPI rollouts are
faster on CPU for this small LSTM on our test machine, so the
deployed controller uses CPU inference.  The numbers in
Table~\ref{tab:mppi_runtime} are measured with a single PyTorch
intra-op thread, which is conservative and leaves CPU headroom for
sensor I/O and IPC.  Using four threads on the same hardware drops
the per-step compute time to 20.07\,ms (about 2\,$\times$ faster;
std 1.77\,ms vs 1.47\,ms) with a comparable overrun rate
($\sim$0.04\%).  We report the single-thread numbers as a conservative
reference; the released evaluator defaults to four threads.  This demonstrates
that learned planning methods can meet stringent real-time
constraints on a commodity host CPU without requiring a GPU at
runtime~\cite{wang2024deployable}.  We caution that the CPU-vs-GPU
ranking reflects our single hardware configuration and may differ on
other machines.

\begin{table}[!ht]
\centering
\caption{MPPI runtime performance (50 hardware trials, single PyTorch
  intra-op thread on the deployment-host CPU, \texttt{calib196\_160} world
  model; the four-thread figure is given in the text). Timings on a different
  CPU will scale accordingly.}
\label{tab:mppi_runtime}
\begin{tabular}{ll}
\toprule
\textbf{Metric} & \textbf{Value} \\
\midrule
Mean compute time & 39.99\,ms \\
Max compute time & 69.81\,ms (cold-start) \\
Median compute time & 39.77\,ms \\
p95 compute time & 42.45\,ms \\
p99 compute time & 43.94\,ms \\
Mean active iteration time & 40.06\,ms \\
Control deadline & 50\,ms (20\,Hz) \\
Device & CPU (PyTorch) \\
\bottomrule
\end{tabular}
\end{table}

Beyond raw speed, MPPI settles cleanly across the difficulty range:
Fig.~\ref{fig:mppi_traj_difficulty} shows its ball-angle trajectories driving
into the $\pm0.01$\,rad band from every stratum, with median settling ranging
from about 1.2\,s (at-wall-same and center starts) to 2.7\,s
(at-wall-opposite), and no sustained oscillation.

\begin{figure}[!t]
  \centering
  \includegraphics[width=0.95\columnwidth]{figures/mppi_trajectories_by_difficulty.eps}
  \caption{MPPI ball-angle trajectories by stratum. Each panel shows all
    trials in the stratum (gray) with the median-settling trial highlighted
    (blue) and the median settling time marked (dotted). The controller
    settles the ball within the $\pm 0.01$\,rad band (green region) across
    all strata, with median settling from about 1.2\,s (at-wall-same and
    center) to 2.7\,s (at-wall-opposite).}
  \label{fig:mppi_traj_difficulty}
\end{figure}
